% Part: first-order-logic
% Chapter: sequent-calculus
% Section: proving-things-quant

\documentclass[../../../include/open-logic-section]{subfiles}

\begin{document}

\olfileid[fr]{fol}{seq}{prq}

\olsection{\usetoken{P}{derivation} avec quantificateurs}

\begin{ex}
Construire une !!{derivation} dans $\Log{LK}$ du séquent
$\lexists[x][\lnot !A(x)] \Sequent \lnot \lforall[x][!A(x)]$.

Avec les quantificateurs, il faut veiller à respecter la condition
sur l'eigenvariable, ce qui peut demander d'ajuster l'ordre de certaines
inférences. En général, il est utile de traiter d'abord, dans la
recherche, les règles soumises à cette condition : elles figureront
plus bas dans la preuve achevée. Il est également utile d'anticiper
la forme du séquent initial. Ici, celui-ci devra être de la forme
$!A(a) \Sequent !A(a)$. Lorsque nous appliquerons les règles pour
les quantificateurs à rebours, nous devrons donc choisir le même
terme, que nous noterons $a$, pour les règles de $\lforall$ et
de $\lexists$. Avec deux termes différents, nous aboutirions à un
séquent tel que $!A(a) \Sequent !A(b)$, qui n'est pas dérivable en
général.\footnote{Note de l'édition française : « en général » précise
la portée du schéma ; certaines instances particulières sont dérivables.}

Comme d'habitude, commençons par écrire :
\begin{prooftree}
\AxiomC{}
\UnaryInf$\lexists[x][\lnot !A(x)] \fCenter \lnot \lforall[x][!A(x)]$
\end{prooftree}
Nous pouvons appliquer soit \LeftR{\exists}, soit \RightR{\lnot}.
Comme \LeftR{\exists} est soumise à la condition sur l'eigenvariable,
traitons-la d'abord.
\begin{prooftree}
\AxiomC{}
\UnaryInf$ \lnot !A(a) \fCenter \lnot \lforall[x][!A(x)]$
\RightLabel{\LeftR{\lexists}}
\UnaryInf$ \lexists[x][\lnot !A(x)] \fCenter \lnot \lforall[x][!A(x)]$
\end{prooftree}
En appliquant à rebours les règles \LeftR{\lnot} et \RightR{\lnot},
nous obtenons :
\begin{prooftree}
\AxiomC{}
\UnaryInf$\lforall[x][!A(x)] \fCenter !A(a)$
\RightLabel{\LeftR{\lnot}}
\UnaryInf$\lnot !A(a), \lforall[x][!A(x)] \fCenter $
\RightLabel{\LeftR{\Exchange}}
\UnaryInf$\lforall[x][!A(x)], \lnot !A(a) \fCenter $
\RightLabel{\RightR{\lnot}}
\UnaryInf$ \lnot !A(a) \fCenter \lnot \lforall[x] !A(x)$
\RightLabel{\LeftR{\lexists}}
\UnaryInf$ \lexists[x] \lnot !A(x) \fCenter \lnot \lforall[x] !A(x)$
\end{prooftree}
Il nous reste à appliquer \LeftR{\forall}. Cette règle n'est pas
soumise à la condition sur l'eigenvariable, et nous pouvons donc
choisir le terme voulu. Puisque nous cherchons un séquent initial
de la forme $!A(a) \Sequent !A(a)$, choisissons $a$ comme argument
de $!A$ lors de cette application.
\begin{prooftree}
\Axiom$!A(a) \fCenter !A(a)$
\RightLabel{\LeftR{\lforall}}
\UnaryInf$\lforall[x][!A(x)] \fCenter !A(a)$
\RightLabel{\LeftR{\lnot}}
\UnaryInf$\lnot !A(a), \lforall[x][!A(x)] \fCenter $
\RightLabel{\LeftR{\Exchange}}
\UnaryInf$\lforall[x][!A(x)], \lnot !A(a) \fCenter $
\RightLabel{\RightR{\lnot}}
\UnaryInf$ \lnot !A(a) \fCenter \lnot \lforall[x][!A(x)]$
\RightLabel{\LeftR{\lexists}}
\UnaryInf$ \lexists[x][ \lnot !A(x)] \fCenter \lnot \lforall[x][!A(x)]$
\end{prooftree}
Vérifions maintenant que la condition sur l'eigenvariable a bien été
respectée, ce qui est essentiel en présence de quantificateurs.
La seule règle utilisée qui soit soumise à cette condition est
\LeftR{\exists}. Or son eigenvariable $a$ ne figure pas dans le
séquent inférieur de cette inférence, qui est ici le séquent final.
La !!{derivation} est donc correcte.
\end{ex}

\begin{prob}
Construire des !!{derivation}s des séquents suivants :
\begin{enumerate}
\item $\Sequent (\lforall[x][!A(x)] \land \lforall[y][!B(y)]) \lif
\lforall[z][(!A(z) \land !B(z))]$.
\item $\Sequent (\lexists[x][!A(x)] \lor \lexists[y][!B(y)]) \lif
\lexists[z][(!A(z) \lor !B(z))]$.
\item $\lforall[x][(!A(x) \lif !B)] \Sequent \lexists[y][!A(y)] \lif !B$.
\item $\lforall[x][\lnot !A(x)] \Sequent \lnot\lexists[x][!A(x)]$.
\item $\Sequent \lnot\lexists[x][!A(x)] \lif \lforall[x][\lnot !A(x)]$.
\item $\Sequent \lnot\lexists[x][\lforall[y][((!A(x,y) \lif \lnot
!A(y,y)) \land (\lnot !A(y,y) \lif !A(x,y)))]]$.
\end{enumerate}
\end{prob}

\begin{prob}
Construire des !!{derivation}s des séquents suivants :
\begin{enumerate}
\item $\Sequent \lnot\lforall[x][!A(x)] \lif \lexists[x][\lnot!A(x)]$.
\item $(\lforall[x][!A(x)] \lif !B) \Sequent \lexists[y][(!A(y) \lif !B)]$.
\item $\Sequent \lexists[x][(!A(x) \lif \lforall[y][!A(y)])]$.
\end{enumerate}
(Tous requièrent la règle $\RightR{\Contraction}$.)
\end{prob}

\end{document}
